% ====== Limitations and Future Work Section ======
\section{Limitations and Future Work}
\label{sec:limitations}
A major limitation appears from attribute-specific articulation patterns: Certain attributes apply exclusively to individual letters, such as \texttt{Istitala} for (ض) and \texttt{Tikrar} for (ر). Consequently, we expect our model will be unable to capture instances of (ض) without \texttt{Istitala} or (ر) without \texttt{Tikrar}. This limitation similarly applies to Tajweed rules that occur less frequently in the Holy Quran, such as \texttt{Imala}, \texttt{Rawm}, and \texttt{Tasheel}. The solution is to annotate real data with these errors.

Additionally, the training dataset comprises only expert male reciters, which may limit generalizability to diverse vocal characteristics. The current scope is restricted to Hafs recitation style; future work will extend to other Riwayah (e.g., Warsh). Finally, qdat\_bench uses a single annotator and one verse; future versions should include multiple experts and broader Quranic coverage.

